% !Mode:: "TeX:UTF-8"
% ============================================================
% INTRODUCTION
% ============================================================
\chapter*{Introduction}
\addcontentsline{toc}{chapter}{Introduction}
\markboth{INTRODUCTION}{INTRODUCTION}

Linear algebra is one of the most widely used branches of modern
mathematics. Matrices, determinants, systems of linear equations, vector
spaces, linear operators and quadratic forms are the working language of
engineering computation, computer graphics and machine learning, and they
are equally basic to \textbf{information security}. The same methods appear
in the classical Hill cipher and in modern lattice-based (post-quantum)
cryptography, in error-detecting and error-correcting codes, and in the
analysis of network traffic and the detection of anomalies.

This study guide accompanies the course “Linear Algebra” and covers the
24 topics of its syllabus. The topics are grouped into five chapters.
Chapter~1, \emph{Matrices and Determinants} (Sections 1.1–1.6), deals with
operations on matrices, determinants of the second, third and $n$-th order,
minors and cofactors, the rank of a matrix and the inverse matrix.
Chapter~2, \emph{Systems of Linear Algebraic Equations, Vectors and Lines in
the Plane} (Sections 2.1–2.6), presents Cramer’s rule, the matrix method and
the Gauss and Gauss–Jordan methods, then turns to vectors, their scalar,
vector and triple products, and the equations of a straight line in the
plane. Chapter~3, \emph{Lines and Planes in Space; Linear and Euclidean
Spaces} (Sections 3.1–3.5), studies lines and planes in space, bases and
subspaces of a linear space, the norm in a Euclidean space and the
construction of orthonormal bases. Chapter~4, \emph{Linear Operators}
(Sections 4.1–4.3), treats linear operators and their matrices, eigenvalues
and eigenvectors, and diagonalization. Chapter~5, \emph{Quadratic Forms,
Second-Order Curves and Iterative Methods} (Sections 5.1–5.4), reduces
quadratic forms and second-order curves to canonical form and closes with
two approximate methods for linear systems, simple (Jacobi) iteration and
the Gauss–Seidel method.

All sections are built in the same way. The main definitions and theorems
come first, followed by worked examples solved in full detail. A separate
subsection shows how the material of the section is used in
\textbf{information security}, so that the reader sees why an engineer in
this field needs it. Each section also contains short examples of the same
computations in Python, Maple and Mathcad. The section ends with a
\textbf{collection of exercises}; every exercise set has 30 variants, which
allows a teacher to give each student in a group an individual assignment,
and the sets are supplemented by applied problems. The answers to all
exercises are collected at the end of the book in “Answers to the
Exercises”, arranged by section and exercise set.

We recommend reading the theory and the worked examples carefully before
attempting the exercises, and comparing one’s own results with the answers
only after the work is done. Checking the computations on a computer, in
Python, MATLAB, Maple or Mathcad, is also useful.

The guide is written for students of information security, computer
engineering, software engineering and other engineering programmes at
universities, and for anyone who studies linear algebra independently.
